%%%%%%%%%%%%%%%%%%%%%%%%%%%%%%%%%%%%%%%%%%%%%%%%%%%%%%%%%%%%%%%%%%%%%%
%%  Copyright by Wenliang Du and SEED Lab contributors.             %%
%%  This work is licensed under the Creative Commons                %%
%%  Attribution-NonCommercial-ShareAlike 4.0 International License. %%
%%  To view a copy of this license, visit                           %%
%%  http://creativecommons.org/licenses/by-nc-sa/4.0/.              %%
%%%%%%%%%%%%%%%%%%%%%%%%%%%%%%%%%%%%%%%%%%%%%%%%%%%%%%%%%%%%%%%%%%%%%%

\newcommand{\commonfolder}{../../../common-files}

\input{\commonfolder/header}
\input{\commonfolder/copyright}
\hypersetup{pdfborder = {0 0 0}}

\lhead{\bfseries SEED Labs -- BGP Prefix Hijacking Lab}

\begin{document}

\begin{center}
{\LARGE BGP Prefix Hijacking Lab}
\end{center}

\seedlabcopyright{2026}


% *******************************************
% SECTION
% *******************************************
\section{Overview}

Border Gateway Protocol (BGP) is the routing protocol used to exchange
reachability information among autonomous systems (ASes) on the Internet. It is
essential to global connectivity, but its trust model is weak: a BGP router can
announce reachability to an IP prefix, and other routers may accept and
propagate that route based on policy rather than cryptographic proof of
ownership.

BGP prefix hijacking happens when an AS announces a prefix that it does not
own. If other networks select the false route, traffic for the victim prefix can
be blackholed, intercepted, or redirected. Many real incidents are caused by
mistakes rather than malicious intent, but the technical effect is the same:
traffic follows an unauthorized route.

This lab is a standalone BGP security lab based on the SEED Internet Emulator.
Students will first study normal route origination, then launch an equal-prefix
hijack and a more-specific prefix hijack from an attacker AS. They will observe
the effect using BIRD route tables, traceroute, ping, and the emulator map.
Finally, students will implement defensive filtering at an upstream provider
and discuss how prefix filters, monitoring, and RPKI/ROV help reduce the risk.

This lab covers the following topics:

\begin{itemize}[noitemsep]
\item BGP route origination and AS paths
\item BGP route selection and longest-prefix matching
\item Equal-prefix and more-specific prefix hijacking
\item Blackholing, redirection, and interception
\item BIRD static routes and import/export filters
\item Defensive filtering and incident response
\end{itemize}

\paragraph{Readings.}
Students should read the SEED BGP tutorial before working on this lab. The
following standards are useful background references:

\begin{itemize}
\item RFC 4271, A Border Gateway Protocol 4 (BGP-4):
      \url{https://www.rfc-editor.org/rfc/rfc4271}
\item RFC 6480, An Infrastructure to Support Secure Internet Routing:
      \url{https://www.rfc-editor.org/rfc/rfc6480}
\item RFC 6811, BGP Prefix Origin Validation:
      \url{https://www.rfc-editor.org/rfc/rfc6811}
\end{itemize}

\paragraph{Lab environment.}
\seedenvironmentB
\nodependency


% *******************************************
% SECTION
% *******************************************
\section{Lab Environment}

This lab is conducted inside the SEED Internet Emulator. The emulator creates a
miniature Internet using Docker containers. Each host and router is a container,
and BGP routers run BIRD.

\paragraph{Starting the emulator.}
Download the BGP emulator lab setup file from the lab website and unzip it. The
actual container files are in the generated \texttt{output} folder. Start the
emulator using Docker Compose:

\begin{lstlisting}
$ cd output
$ docker-compose build
$ docker-compose up

// Aliases in the SEED VM
$ dcbuild
$ dcup
$ dcdown
\end{lstlisting}

After the emulator starts, the network map can be accessed from:

\begin{lstlisting}
http://localhost:8080/map.html
\end{lstlisting}

\paragraph{Emulator conventions.}
The emulator uses several naming and addressing conventions:

\begin{itemize}[noitemsep]
\item AS numbers \texttt{2}--\texttt{9}: large transit ASes.
\item AS numbers \texttt{10}--\texttt{19}: smaller transit ASes.
\item AS numbers \texttt{100}--\texttt{149}: Internet Exchanges.
\item AS numbers \texttt{150}--\texttt{199}: stub ASes.
\item The first prefix for AS-\texttt{N} is \texttt{10.N.0.0/24}.
\item In each AS, host addresses often start at \texttt{71}; router addresses
      are usually near the end of the subnet.
\end{itemize}

In this lab, AS-154 is the victim, and AS-161 is the attacker. The victim
prefix is \texttt{10.154.0.0/24}. A typical victim host is
\texttt{10.154.0.71}.

\paragraph{Useful commands.}
The following commands are used throughout the lab:

\begin{lstlisting}
// List containers
$ dockps

// Get a shell inside a container
$ docksh <container-id>

// Show BGP sessions
# birdc show protocols

// Show a route
# birdc show route all for 10.154.0.0/24

// Reload BIRD configuration
# birdc configure

// Test reachability and path
$ ping 10.154.0.71
$ traceroute 10.154.0.71
\end{lstlisting}

\paragraph{Modifying BIRD configuration.}
Students can edit \path{/etc/bird/bird.conf} directly inside a container, or
copy it to the host, edit it, and copy it back:

\begin{lstlisting}
$ docker cp <container-id>:/etc/bird/bird.conf ./bird.conf

... edit bird.conf ...

$ docker cp ./bird.conf <container-id>:/etc/bird/bird.conf
$ docker exec <container-id> birdc configure
\end{lstlisting}

\paragraph{Helper snippets.}
This lab folder contains a \texttt{Labsetup} folder with small BIRD snippets:

\begin{lstlisting}
Labsetup/
|-- as161-hijack-10.154.0.0-24.conf
|-- as161-hijack-10.154.0.0-25.conf
|-- as3-filter-as154-prefix.conf
\end{lstlisting}

These snippets are not complete configuration files. They are examples to paste
into the appropriate part of the existing BIRD configuration.


% *******************************************
% SECTION
% *******************************************
\section{Theory: How BGP Prefix Hijacking Works}

\subsection{Route Origination}

An AS originates a route when it announces that it can reach a prefix. For
example, AS-154 normally originates \texttt{10.154.0.0/24}. When other BGP
routers receive this announcement, they store the route together with path
attributes, including the AS path.

The AS path records the sequence of ASes that the announcement has traversed.
The last AS in the path is the origin AS. If a router sees a path such as
\texttt{3 154}, the origin AS is AS-154. If it sees \texttt{3 161}, the origin
AS is AS-161.

Traditional BGP does not prove that the origin AS owns the prefix. This is the
root of prefix hijacking.

\subsection{Route Selection}

When a router learns multiple routes to the same prefix, it runs the BGP route
selection process. The complete process has many steps, but common factors
include:

\begin{itemize}[noitemsep]
\item Local preference
\item AS path length
\item Origin type
\item MED
\item eBGP over iBGP
\item IGP cost to the next hop
\end{itemize}

The most important point for this lab is that BGP may select an unauthorized
route if that route is more preferred by policy. Shorter AS paths often help an
attacker, but local policy can override path length.

\subsection{Longest-Prefix Matching}

BGP route selection chooses the best route for a given prefix. Packet forwarding
uses another rule: longest-prefix matching. If a router has routes for both
\texttt{10.154.0.0/24} and \texttt{10.154.0.0/25}, a packet to
\texttt{10.154.0.71} matches both, but the \texttt{/25} route is more specific
and will be used.

This makes more-specific hijacking powerful. Even if the victim's
\texttt{/24} route remains visible, an attacker who successfully announces a
covered \texttt{/25} may capture traffic for half of the victim's address
space.

\subsection{Types of Prefix Hijacking}

This lab studies two common forms:

\begin{itemize}[noitemsep]
\item \textbf{Equal-prefix hijack:} the attacker announces the same prefix as
      the victim, such as \texttt{10.154.0.0/24}.
\item \textbf{More-specific hijack:} the attacker announces a more-specific
      prefix, such as \texttt{10.154.0.0/25}.
\end{itemize}

Equal-prefix hijacking depends heavily on BGP route selection. More-specific
hijacking depends heavily on longest-prefix matching and can be more disruptive.

\subsection{Attack Outcomes}

A prefix hijack can have several outcomes:

\begin{itemize}[noitemsep]
\item \textbf{Blackholing:} traffic is routed to the attacker and dropped.
\item \textbf{Interception:} traffic is routed through the attacker and then
      forwarded to the victim, making the attack harder to notice.
\item \textbf{Impersonation:} the attacker hosts a fake service for the victim
      prefix.
\item \textbf{Partial outage:} only some ASes select the false route.
\end{itemize}

In this lab, the basic attack uses blackhole routes because they are simple and
safe to demonstrate inside the emulator. Instructors can extend the lab to
interception by making the attacker forward traffic onward to the victim.

\subsection{Defenses}

Several defenses can reduce prefix hijacking:

\begin{itemize}[noitemsep]
\item Prefix filters at customer/provider boundaries.
\item Monitoring for unexpected origin ASes or sudden path changes.
\item RPKI and Route Origin Validation (ROV), which can reject routes whose
      origin AS is not authorized by a ROA.
\item Incident response procedures, including contacting upstream providers.
\end{itemize}

This lab focuses on prefix filters because they are easy to implement directly
in BIRD. RPKI/ROV is discussed as a stronger ecosystem-level defense and can be
studied in a separate lab.


% *******************************************
% SECTION
% *******************************************
\section{Lab Tasks}

% -------------------------------------------
% SUBSECTION
% -------------------------------------------
\subsection{Task 1: Observe Normal Routing to the Victim Prefix}

The objective of this task is to establish a baseline before the attack.
Choose a host outside AS-154, such as a host in AS-155 or AS-156, and test
connectivity to the victim host:

\begin{lstlisting}
$ ping 10.154.0.71
$ traceroute 10.154.0.71
\end{lstlisting}

Then inspect the route to AS-154's prefix on a nearby BGP router:

\begin{lstlisting}
# birdc show route all for 10.154.0.0/24
\end{lstlisting}

Please answer the following questions:

\begin{itemize}[noitemsep]
\item Which AS originates \texttt{10.154.0.0/24}?
\item What AS path do you observe?
\item Which next hop is used?
\item Does traffic reach \texttt{10.154.0.71} successfully?
\end{itemize}

\paragraph{Connection to theory.}
Identify the origin AS using the last AS in the AS path. Explain why this
baseline route is expected to be legitimate.


% -------------------------------------------
% SUBSECTION
% -------------------------------------------
\subsection{Task 2: Launch an Equal-Prefix Hijack from AS-161}

In this task, AS-161 announces the victim's prefix
\texttt{10.154.0.0/24}. This is an equal-prefix hijack.

Get a shell on AS-161's BGP router and add the following block to
\path{/etc/bird/bird.conf}:

\begin{lstlisting}
protocol static {
    ipv4 {
        table t_bgp;
    };
    route 10.154.0.0/24 blackhole {
        bgp_large_community.add(LOCAL_COMM);
    };
}
\end{lstlisting}

Reload BIRD:

\begin{lstlisting}
# birdc configure
\end{lstlisting}

Wait for the route to propagate. From at least two ASes outside AS-154, inspect
the route and test connectivity:

\begin{lstlisting}
# birdc show route all for 10.154.0.0/24
$ traceroute 10.154.0.71
$ ping 10.154.0.71
\end{lstlisting}

Please provide screenshots of the route table and traffic tests. If some ASes
select the false route and others do not, record the difference.

\paragraph{Connection to theory.}
Explain how BGP route selection decides between the legitimate route from
AS-154 and the false route from AS-161. Do not simply say that the attacker
"wins"; identify the observed AS path and next hop.


% -------------------------------------------
% SUBSECTION
% -------------------------------------------
\subsection{Task 3: Visualize and Capture the Hijack}

The objective of this task is to collect evidence of the hijack. Use the
emulator map and packet capture facilities to visualize traffic to AS-154.
Set a filter such as:

\begin{lstlisting}
icmp
\end{lstlisting}

Generate traffic from a host outside AS-154:

\begin{lstlisting}
$ ping 10.154.0.71
\end{lstlisting}

Please describe where the packets go before and after the hijack. If you use
\texttt{tcpdump}, save a short packet capture and identify which routers see
the traffic.

\paragraph{Connection to theory.}
Connect the packet path to the selected BGP route. Explain why packet
forwarding changes after the control-plane route changes.


% -------------------------------------------
% SUBSECTION
% -------------------------------------------
\subsection{Task 4: Launch a More-Specific Prefix Hijack}

In this task, AS-161 announces \texttt{10.154.0.0/25}. This is a
more-specific prefix contained inside AS-154's \texttt{/24}.

On AS-161's BGP router, replace the previous malicious static route with:

\begin{lstlisting}
protocol static {
    ipv4 {
        table t_bgp;
    };
    route 10.154.0.0/25 blackhole {
        bgp_large_community.add(LOCAL_COMM);
    };
}
\end{lstlisting}

Reload BIRD and observe the routes:

\begin{lstlisting}
# birdc configure
# birdc show route all for 10.154.0.0/24
# birdc show route all for 10.154.0.0/25
\end{lstlisting}

Test two addresses:

\begin{lstlisting}
$ traceroute 10.154.0.71
$ traceroute 10.154.0.171
\end{lstlisting}

The address \texttt{10.154.0.71} is inside \texttt{10.154.0.0/25}, while
\texttt{10.154.0.171} is outside that \texttt{/25} but still inside the
\texttt{/24}. If the emulator does not have a host at one of these addresses,
use route-table evidence and choose another address that demonstrates the same
prefix-match property.

\paragraph{Connection to theory.}
Explain why a more-specific \texttt{/25} can affect forwarding even when the
legitimate \texttt{/24} route still exists. Your answer should explicitly use
the term longest-prefix matching.


% -------------------------------------------
% SUBSECTION
% -------------------------------------------
\subsection{Task 5: Victim Response}

AS-154 detects that its traffic is being diverted. Assume the operators of
AS-154 cannot immediately reach AS-161's provider. They decide to try to
restore traffic by changing their own announcements.

Design and test one response from AS-154. Examples include:

\begin{itemize}[noitemsep]
\item Announce a more-specific prefix that covers the affected address range.
\item Change export policy so the victim route is propagated through another
      provider if one exists.
\item Use the emulator map to identify which ASes still select the legitimate
      route and which select the attack route.
\end{itemize}

Please describe what you changed, what improved, and what did not improve.

\paragraph{Connection to theory.}
Explain why announcing more-specific prefixes can help the victim recover
traffic, and explain why this is not a good long-term defense for the Internet.


% -------------------------------------------
% SUBSECTION
% -------------------------------------------
\subsection{Task 6: Provider-Side Filtering Defense}

Eventually, AS-3, the upstream provider for AS-161, is contacted. AS-3 does not
want to disconnect AS-161 completely, but it wants to stop AS-161 from
announcing AS-154's prefix.

On AS-3's BGP router, identify the BGP session with AS-161:

\begin{lstlisting}
# birdc show protocols
\end{lstlisting}

Add a filter condition to the import filter for AS-161's session:

\begin{lstlisting}
if net ~ [ 10.154.0.0/24{24,32} ] then {
    print "Rejecting unauthorized AS-154 prefix from AS-161: ", net;
    reject;
}
\end{lstlisting}

Reload BIRD:

\begin{lstlisting}
# birdc configure
\end{lstlisting}

Repeat the attacks from Tasks 2 and 4. Please answer the following:

\begin{itemize}[noitemsep]
\item Does AS-3 reject the false route?
\item Does AS-161 still have connectivity for its own prefix
      \texttt{10.161.0.0/24}?
\item Does traffic to AS-154 return to the legitimate route?
\item What evidence from \texttt{birdc} or traffic tests supports your answer?
\end{itemize}

\paragraph{Connection to theory.}
Explain why the filter is a policy defense at the provider boundary. Which
announcements are blocked, and which legitimate routes are still allowed?


% -------------------------------------------
% SUBSECTION
% -------------------------------------------
\subsection{Task 7: Monitoring and Incident Report}

In this task, students act as network operators. Create a short incident report
for the hijack. The report should include:

\begin{itemize}[noitemsep]
\item The victim prefix.
\item The legitimate origin AS.
\item The unauthorized origin AS.
\item The time or sequence of observations.
\item Evidence from route tables, traceroute, ping, packet capture, or the map.
\item The mitigation applied.
\item Residual risks or limitations.
\end{itemize}

\paragraph{Connection to theory.}
Explain how the evidence distinguishes a BGP control-plane problem from a host
failure or ordinary packet loss.


% *******************************************
% SECTION
% *******************************************
\section{Guidelines}

\paragraph{Removing the attack route.}
After finishing an attack task, remove or comment out the malicious
\texttt{protocol static} block and reload BIRD:

\begin{lstlisting}
# birdc configure
\end{lstlisting}

You can also restart the emulator to return to a clean state.

\paragraph{Finding the right router.}
Use \texttt{dockps} and the emulator map to identify the router container for
an AS. Router container names usually contain the AS number and the word
\texttt{router}.

\paragraph{Understanding \texttt{LOCAL\_COMM}.}
In the SEED emulator, some export filters use large communities to decide which
routes should be exported. If a static route is not added to
\texttt{LOCAL\_COMM}, it may stay local and may not be exported to BGP peers.
This is why the malicious static route includes:

\begin{lstlisting}
bgp_large_community.add(LOCAL_COMM);
\end{lstlisting}

\paragraph{Equal-prefix versus more-specific hijacking.}
If your equal-prefix hijack does not affect every AS, this is not necessarily a
mistake. Different ASes may select different routes depending on policy and AS
path. A more-specific hijack usually has broader forwarding impact because of
longest-prefix matching.

\paragraph{Instructor customization.}
Instructors can change the victim AS, attacker AS, provider AS, and prefixes.
For example, AS-155 can be used as the victim and AS-162 as the attacker. If
the AS numbers change, update the prefix values using the emulator convention:
AS-\texttt{N}'s first prefix is \texttt{10.N.0.0/24}.


% *******************************************
% SECTION
% *******************************************
\section{Submission}

%%%%%%%%%%%%%%%%%%%%%%%%%%%%%%%%%%%%%%%%
\input{\commonfolder/submission}
%%%%%%%%%%%%%%%%%%%%%%%%%%%%%%%%%%%%%%%%

In addition, your report should include:

\begin{itemize}[noitemsep]
\item Baseline route and traffic evidence before the attack.
\item Evidence for the equal-prefix hijack.
\item Evidence for the more-specific hijack.
\item The BIRD configuration changes you made.
\item Evidence that provider-side filtering stops the false announcement.
\item A short incident report for the attack.
\end{itemize}

\end{document}
